\chapter{Introduzione}
\label{chap:introduzione}

Il settore marittimo rappresenta l'infrastruttura portante del commercio globale, ma si trova oggi ad affrontare una fase di profonda trasformazione tecnologica orientata all'efficienza, alla sicurezza e alla sostenibilità. In questo contesto, l'emergere delle \ac{MASS} promette di rivoluzionare la navigazione riducendo l'errore umano e ottimizzando i processi operativi~\cite{11217861, roadmap}. Tuttavia, l'inadeguatezza delle reti di comunicazione marittima convenzionali frena il passaggio verso navi interamente o parzialmente autonome. Mentre i moderni dispositivi a bordo (come telecamere 4K, \ac{LiDAR} e radar) generano massicci flussi di dati necessari per mantenere un'adeguata situational awareness, le comunicazioni tradizionali offrono una larghezza di banda limitata, insufficiente a supportare lo scambio di streaming video o di point clouds in tempo reale~\cite{11217861, 10.1145/2664243.2664257}. Per superare questo collo di bottiglia senza dipendere da costose connessioni satellitari, emerge la necessità di adottare paradigmi come Aggregate Programming, puntando a creare reti multi-hop e auto-organizzanti direttamente tra le imbarcazioni e le infrastrutture costiere~\cite{beal2015aggregate, casadei2019self}. 

La motivazione alla base di questo lavoro di tesi nasce dall'analisi dei limiti delle attuali architetture di rete in ambito marittimo. Poiché in questo tipo di ambiente la connettività è sparsa e i data rate sono bassi, l'articolo \emph{Robust Communication through Collective Adaptive Relay Schemes for Maritime Vessels} ha valutato approcci in cui i dati vengono raccolti localmente e poi riassunti in punti strategici della rete~\cite{11217861}. Idealmente, le imbarcazioni si auto-organizzano in gruppi (definiti \emph{cluster}) all'interno dei quali i dati possono essere trasmessi in modo affidabile verso un nodo designato come sintetizzatore, ovvero il leader del cluster~\cite{11217861}. Il leader ha il compito di comprimere i flussi raccolti, incluso il proprio, e inoltrare questo output sintetizzato alla stazione a terra o a un relay nel cluster successivo, sfruttando torri cellulari 5G~\cite{11217861}. Sebbene questa architettura (identificata come la ``Fase 1" dell'agenda \emph{CoMPass}) migliori il data rate in scenari simulati, presenta una debolezza: la formazione dei cluster è basata solo su posizioni statiche (che non tengono conto della velocità o della rotta delle navi)~\cite{11217861, roadmap}. L'uso di queste metriche non è infatti sufficiente, in quanto le imbarcazioni con traiettorie e velocità differenti dovrebbero essere scoraggiate dal formare collegamenti di comunicazione che diventerebbero rapidamente instabili~\cite{roadmap}. 

L'obiettivo principale di questo lavoro di tesi è implementare la ``Fase 2" della roadmap CoMPass, ossia la riconfigurazione della rete consapevole della navigazione~\cite{roadmap}. Partendo dal framework di comunicazione stabilito, questa fase incorpora il posizionamento dinamico delle navi con l'intento di rendere più realistica e selettiva la topologia della rete. Nello specifico, questa modifica utilizza i dati di navigazione estraendo dal payload dei messaggi \ac{AIS} i valori di \ac{SOG} e \ac{COG} per perfezionare gli algoritmi di clustering e le strategie di comunicazione. Valutando questi parametri, le navi che si muovono in direzioni e velocità simili ottengono la priorità per l'appartenenza al cluster mentre le imbarcazioni con traiettorie e velocità divergenti vengono scoraggiate dal formare collegamenti di comunicazione che diventerebbero rapidamente obsoleti.

Il documento è strutturato in capitoli secondo la seguente organizzazione. Il \cref{chap:background} espone i fondamenti teorici e le tecnologie di riferimento, descrivendo i sistemi di comunicazione marittima convenzionali (con particolare focus sulla struttura e i messaggi \ac{AIS}), Aggregate Programming, Field Calculus e gli strumenti di simulazione utilizzati. Viene inoltre descritto in dettaglio l'esperimento precedente. Il \cref{chap:analisi} è dedicato all'analisi del problema. Vengono descritti: i limiti dell'algoritmo esistente, gli obiettivi della tesi e i requisiti funzionali e non funzionali richiesti per la realizzazione del nuovo sistema. Il \cref{chap:designImplementazione} espone il design e l'implementazione, illustrando la metodologia di utilizzo e integrazione dei nuovi dati all'interno del programma e le modifiche all'algoritmo. Il \cref{chap:valutazioni} presenta la comparazione dei risultati ottenuti con quelli precedenti, utilizzando le stesse metriche di valutazione. Infine si discute delle conclusioni e dei possibili sviluppi futuri del lavoro svolto.
